\section{The Appointed Texts}
\label{sec:appointed-texts}

The Missal's texts for this Sunday are those of \latin{Dominica XXVIII ``per
annum''} in the \work{Missale Romanum}, editio typica tertia (2002), which the
\work{Roman Missal}, Third Edition, renders for the dioceses of the United
States; the readings, psalm and Alleluia verse are those of no.~142 of the
\work{Lectionary for Mass}, which follows the \work{Ordo lectionum Missae} of
1981. The Latin of the Missal, the approved English of the Missal and the
approved English of the Lectionary are all under copyright, and none of them
is reproduced here. Each biblical element is given at its canonical verses in
the Douay--Rheims version as revised by Bishop Challoner, which translates the
Clementine Vulgate and numbers the psalms as the Vulgate does. That English is
a study text and is not the wording proclaimed at Mass. Words printed in
brackets and italics belong to the verse but not to the appointed text. Each
oration is identified by the opening words of its Latin and described, and no
English is supplied for it; the three antiphons and the Alleluia verse adapt
Scripture, and the verses they draw on are printed beside them as their
sources and not as their words.

\Needspace{8\baselineskip}
\subsection{Entrance Antiphon}
\label{proper-entrance-antiphon}

\properlocus{\latin{Si iniquitates observaveris, Domine}. The Missal cites
Psalm 129:3--4 in the Vulgate numbering, which is Psalm 130:3--4 in the
Hebrew, and prints no \latin{Cf.}}

\noindent A question put to the Lord: if he were to keep account of
iniquities, who could stand? The answer follows, that with him there is
propitiation, and the antiphon closes by addressing him as the God of Israel.
It takes v.~3 of the psalm and the first clause of v.~4, and then departs from
the psalm: it leaves out the rest of v.~4, on waiting for the Lord by reason of
his law, and ends with the vocative \latin{Deus Israel}, which stands neither in
the Clementine Vulgate's Psalm 129 nor in the \work{Nova Vulgata}'s. Within the
psalm Israel is named only in vv.~6 and~8, in the third person. The antiphon is
therefore an adaptation, though it carries no \latin{Cf.}; the Latin source of
its ending has not been established. No English is printed for the antiphon or
its closing words. The verses it draws on follow.

\begin{englishwitness}{Douay--Rheims (Challoner), Ps 129:3--4, the antiphon's source}
\textsuperscript{3}If thou, O Lord, wilt mark iniquities: Lord, who shall stand
it. \textsuperscript{4}For with thee there is merciful forgiveness\notread{: and
by reason of thy law, I have waited for thee, O Lord. My soul hath relied on
his word:}
\end{englishwitness}

\noindent ``Merciful forgiveness'' is the Douay--Rheims rendering of
\latin{propitiatio}.

\Needspace{8\baselineskip}
\subsection{Collect}
\label{proper-collect}

\properlocus{\latin{Tua nos, quaesumus, Domine, gratia}. Long conclusion.}

\noindent A petition about grace and action: that the Lord's grace always go
before the assembly and follow it (\latin{semper et praeveniat et sequatur}),
and keep it constantly intent on good works (\latin{bonis operibus iugiter
praestet esse intentos}). No English is printed.

\Needspace{8\baselineskip}
\subsection{First Reading}
\label{proper-first-reading}

\properlocus{Isaiah 25:6--10a. Lectionary no.~142. The reading ends with the
first clause of v.~10; the rest of the verse, on Moab, is not read.}

\begin{englishwitness}{Douay--Rheims (Challoner), Is 25:6--10}
\textsuperscript{6}And the Lord of hosts shall make unto all people in this
mountain, a feast of fat things, a feast of wine, of fat things full of marrow,
of wine purified from the lees. \textsuperscript{7}And he shall destroy in this
mountain the face of the bond with which all people were tied, and the web that
he began over all nations. \textsuperscript{8}He shall cast death down headlong
for ever: and the Lord God shall wipe away tears from every face, and the
reproach of his people he shall take away from off the whole earth: for the
Lord hath spoken it. \textsuperscript{9}And they shall say in that day: Lo,
this is our God, we have waited for him, and he will save us: this is the Lord,
we have patiently waited for him, we shall rejoice and be joyful in his
salvation. \textsuperscript{10}For the hand of the Lord shall rest in this
mountain\notread{: and Moab shall be trodden down under him, as straw is broken
in pieces with the wain}.
\end{englishwitness}

\Needspace{8\baselineskip}
\subsection{Responsorial Psalm}
\label{proper-responsorial-psalm}

\properlocus{Psalm 23 (22):1--3a, 3b--4, 5, 6; the response is taken from
v.~6cd. Lectionary no.~142.}

\begin{englishwitness}{Douay--Rheims (Challoner), Ps 22:6, second half, as the response}
\notread{And thy mercy will follow me all the days of my life.} And that I may
dwell in the house of the Lord unto length of days.
\end{englishwitness}

\noindent The Latin \work{Ordo} prints the response from the second half of
v.~6, the dwelling in the house of the Lord for length of days. The response on
the bishops' conference's page of readings for the day gives the same
locator, but its wording joins the dwelling in the Lord's house to the days of
one's life, which in the Latin and in the Douay--Rheims belong to the first
half of the verse. Whether the printed Lectionary reads the same was not
established, and nothing below rests on the wording of the response.

\begin{englishwitness}{Douay--Rheims (Challoner), Ps 22:1--6}
\textsuperscript{1}\notread{A psalm for David.} The Lord ruleth me: and I
shall want nothing. \textsuperscript{2}He hath set me in a place of pasture. He
hath brought me up, on the water of refreshment: \textsuperscript{3}He hath
converted my soul.

He hath led me on the paths of justice, for his own name's sake.
\textsuperscript{4}For though I should walk in the midst of the shadow of
death, I will fear no evils, for thou art with me. Thy rod and thy staff, they
have comforted me.

\textsuperscript{5}Thou hast prepared a table before me against them that
afflict me. Thou hast anointed my head with oil; and my chalice which
inebriateth me, how goodly is it!

\textsuperscript{6}And thy mercy will follow me all the days of my life. And
that I may dwell in the house of the Lord unto length of days.
\end{englishwitness}

\noindent The strophes divide within v.~3, between the converting of the soul
and the leading in the paths of justice. The title, which the Vulgate numbers
within v.~1, is not read. Where the Latin says that the Lord rules the singer,
the Hebrew says that he shepherds him.

\Needspace{8\baselineskip}
\subsection{Second Reading}
\label{proper-second-reading}

\properlocus{Philippians 4:12--14, 19--20. Lectionary no.~142. Verses 15--18
are left out between the two parts of the reading.}

\begin{englishwitness}{Douay--Rheims (Challoner), Phil 4:12--14, 19--20}
\textsuperscript{12}I know both how to be brought low, and I know how to abound
(every where and in all things I am instructed): both to be full and to be
hungry: both to abound and to suffer need. \textsuperscript{13}I can do all
things in him who strengtheneth me. \textsuperscript{14}Nevertheless, you have
done well in communicating to my tribulation.

\notread{Verses 15--18 are not read.}

\textsuperscript{19}And may my God supply all your want, according to his
riches in glory in Christ Jesus. \textsuperscript{20}Now to God and our Father
be glory, world without end. Amen.
\end{englishwitness}

\noindent In v.~19 the Clementine Vulgate, which the Douay--Rheims follows,
has a wish, \latin{impleat omne desiderium vestrum}; the \work{Nova Vulgata},
the Latin of the \work{Ordo}, has a promise, \latin{implebit}, and the
Lectionary's English, as the bishops' conference's page prints it, is likewise
a promise. Where an argument below turns on the mood of the verb or on the
word \latin{desiderium}, it says which text it follows.

\Needspace{8\baselineskip}
\subsection{Alleluia Verse}
\label{proper-gospel-acclamation}

\properlocus{Cf.\ Ephesians 1:17--18. Lectionary no.~142.}

\noindent The verse makes Paul's prayer for his readers the assembly's own
prayer. It names the Father of our Lord Jesus Christ as the one who is to
enlighten the eyes of the heart, leaves out the spirit of wisdom and
revelation, and puts the whole in the first person plural, so that it ends on
the hope of \emph{our} calling (\latin{spes vocationis nostrae}) where Paul
wrote of the hope of \emph{his} calling. No English is printed for the verse.
The verses it adapts follow.

\begin{englishwitness}{Douay--Rheims (Challoner), Eph 1:17--18, the verse's source}
\textsuperscript{17}That the God of our Lord Jesus Christ, the Father of glory,
\notread{may give unto you the spirit of wisdom and of revelation, in the
knowledge of him:} \textsuperscript{18}The eyes of your heart enlightened that
you may know what the hope is of his calling \notread{and what are the riches
of the glory of his inheritance in the saints.}
\end{englishwitness}

\Needspace{8\baselineskip}
\subsection{Gospel, Longer Form}
\label{proper-gospel-long}

\properlocus{Matthew 22:1--14. Lectionary no.~142, the longer form. The Latin
\work{Ordo} opens the reading with an incipit naming the chief priests and
elders of the people, whom the narrative introduces at 21:23; the incipit is
not printed here as Scripture.}

\begin{englishwitness}{Douay--Rheims (Challoner), Mt 22:1--14}
\textsuperscript{1}And Jesus answering, spoke again in parables to them,
saying: \textsuperscript{2}The kingdom of heaven is likened to a king who made
a marriage for his son. \textsuperscript{3}And he sent his servants to call them
that were invited to the marriage: and they would not come.
\textsuperscript{4}Again he sent other servants, saying: Tell them that were
invited, Behold, I have prepared my dinner: my beeves and fatlings are killed,
and all things are ready. Come ye to the marriage. \textsuperscript{5}But they
neglected and went their ways, one to his farm and another to his merchandise.
\textsuperscript{6}And the rest laid hands on his servants and, having treated
them contumeliously, put them to death. \textsuperscript{7}But when the king
had heard of it, he was angry: and sending his armies, he destroyed those
murderers and burnt their city. \textsuperscript{8}Then he saith to his
servants: The marriage indeed is ready; but they that were invited were not
worthy. \textsuperscript{9}Go ye therefore into the highways; and as many as
you shall find, call to the marriage. \textsuperscript{10}And his servants
going forth into the ways, gathered together all that they found, both bad and
good: and the marriage was filled with guests. \textsuperscript{11}And the king
went in to see the guests: and he saw there a man who had not on a wedding
garment. \textsuperscript{12}And he saith to him: Friend, how camest thou in
hither not having on a wedding garment? But he was silent.
\textsuperscript{13}Then the king said to the waiters: Bind his hands and feet,
and cast him into the exterior darkness. There shall be weeping and gnashing of
teeth. \textsuperscript{14}For many are called, but few are chosen.
\end{englishwitness}

\Needspace{8\baselineskip}
\subsection{Gospel, Shorter Form}
\label{proper-gospel-short}

\properlocus{Matthew 22:1--10. Lectionary no.~142, the shorter form, with the
same incipit.}

\begin{englishwitness}{Douay--Rheims (Challoner), Mt 22:1--10}
\textsuperscript{1}And Jesus answering, spoke again in parables to them,
saying: \textsuperscript{2}The kingdom of heaven is likened to a king who made
a marriage for his son. \textsuperscript{3}And he sent his servants to call them
that were invited to the marriage: and they would not come.
\textsuperscript{4}Again he sent other servants, saying: Tell them that were
invited, Behold, I have prepared my dinner: my beeves and fatlings are killed,
and all things are ready. Come ye to the marriage. \textsuperscript{5}But they
neglected and went their ways, one to his farm and another to his merchandise.
\textsuperscript{6}And the rest laid hands on his servants and, having treated
them contumeliously, put them to death. \textsuperscript{7}But when the king
had heard of it, he was angry: and sending his armies, he destroyed those
murderers and burnt their city. \textsuperscript{8}Then he saith to his
servants: The marriage indeed is ready; but they that were invited were not
worthy. \textsuperscript{9}Go ye therefore into the highways; and as many as
you shall find, call to the marriage. \textsuperscript{10}And his servants
going forth into the ways, gathered together all that they found, both bad and
good: and the marriage was filled with guests.
\end{englishwitness}

\noindent The shorter form ends with the hall filled. It does not contain the
guest without a wedding garment or the saying ``many are called, but few are
chosen''.

\Needspace{8\baselineskip}
\subsection{Prayer over the Offerings}
\label{proper-prayer-over-offerings}

\properlocus{\latin{Suscipe, Domine, fidelium preces}. Short conclusion.}

\noindent The Lord is asked to receive the prayers of the faithful together
with the sacrificial offerings, so that through these acts of devout worship
(\latin{per haec piae devotionis officia}) they may pass over to heavenly glory
(\latin{ad caelestem gloriam}). No English is printed.

\Needspace{8\baselineskip}
\subsection{Communion Antiphon, First Option}
\label{proper-communion-antiphon-a}

\properlocus{\latin{Divites eguerunt et esurierunt}. The Missal cites
\latin{Cf.} Psalm 33:11 in the Vulgate numbering, which is Psalm 34:11 in the
Hebrew count that numbers the title, and 34:10 in English Bibles that do not.}

\noindent The rich have known want and hunger, but those who seek the Lord will
lack no good thing. The antiphon reads \latin{quaerentes}, those who seek, where
the Clementine Vulgate has \latin{inquirentes} and the \work{Nova Vulgata} has
\latin{inquirentes} with another verb for lacking; the rest follows the
Clementine. The Hebrew text, as the King James Version translates it, has young
lions where the Latin and the Douay--Rheims have the rich, so that what is said
of the rich here belongs to the Latin and Greek tradition the Missal cites. No
English is printed for the antiphon.

\begin{englishwitness}{Douay--Rheims (Challoner), Ps 33:11, the antiphon's source}
The rich have wanted, and have suffered hunger: but they that seek the Lord
shall not be deprived of any good.
\end{englishwitness}

\Needspace{8\baselineskip}
\subsection{Communion Antiphon, Second Option}
\label{proper-communion-antiphon-b}

\properlocus{\latin{Cum apparuerit Dominus}. 1~John 3:2, printed under
\latin{Vel} and without \latin{Cf.}}

\noindent When the Lord appears, those who are his will be like him, because
they will see him as he is. The antiphon takes only the second half of the
verse, leaves out its opening words ``We know that'', and supplies a subject,
\latin{Dominus}, which neither the Clementine nor the \work{Nova Vulgata}
expresses. It is therefore an adaptation, though it carries no \latin{Cf.} No
English is printed for the antiphon. The verse it adapts follows.

\begin{englishwitness}{Douay--Rheims (Challoner), 1~Jn 3:2, the antiphon's source}
\notread{Dearly beloved, we are now the sons of God: and it hath not yet
appeared what we shall be. We know that} when he shall appear we shall be like
to him: because we shall see him as he is.
\end{englishwitness}

\Needspace{8\baselineskip}
\subsection{Prayer after Communion}
\label{proper-prayer-after-communion}

\properlocus{\latin{Maiestatem tuam, Domine, suppliciter deprecamur}. Short
conclusion.}

\noindent The prayer is addressed to the divine majesty: as he feeds the
assembly with the Body and Blood of Christ, so may he make them partakers of
the divine nature. The Missal prints no source for it; its closing words,
\latin{divinae naturae \ldots\ consortes}, are those of 2~Peter 1:4 in the
Clementine Vulgate, ``that by these you may be made partakers of the divine
nature''. No English is printed for the prayer.
